% Method section expanded for the moderate-length revision.
% This file is a section fragment, not a standalone LaTeX document.

\section{Method}
\label{sec:method}

\subsection{Overview}
\label{subsec:method_overview}

The evaluated pipeline starts from a trained FP32 model and applies simulated quantization to weights and activations. We then evaluate the resulting fake-quantized forward pass. No model weights are updated, no quantization-aware training is performed, and the test set is not used for calibration. The evaluated method changes only the activation threshold selection rule during calibration.

The main idea is to replace a single MinMax activation range with a layer-wise threshold selected from a small candidate set. This design keeps the calibration procedure simple and reproducible. It also keeps the method conservative: the selected threshold is the best candidate only within the tested percentile set, not a global optimum over all possible clipping thresholds.

\subsection{Simulated Quantization}
\label{subsec:quantization_setting}

The INT4 activation experiments quantize post-ReLU activations as unsigned 4-bit values in the range $[0,15]$. Because the observed activation sites are post-ReLU modules, the clipping interval for activation site $l$ is $[0, \alpha_l]$. Given an activation value $a$, the simulated activation quantizer is
\begin{equation}
\begin{aligned}
\tilde{a} &= s_l\cdot\operatorname{clip}\left(
\operatorname{round}\left(\frac{\operatorname{clip}(a,0,\alpha_l)}{s_l}\right),0,15\right), \\
 s_l &= \frac{\alpha_l}{15}.
\end{aligned}
\label{eq:activation_fake_quant}
\end{equation}
The post-ReLU activation zero-point is fixed to zero; equivalently, this is an unsigned affine form specialized to non-negative activations. Weights in convolutional and linear layers are fake-quantized using per-output-channel symmetric quantization. The INT4 weight range is $[-7,7]$, and the INT8 baseline uses the corresponding 8-bit ranges specified in the configuration.

\subsection{Activation Calibration Rules}
\label{subsec:activation_calibration_rules}

For an INT4 post-ReLU activation site $l$, let $\mathcal{A}_l=\{a_{l,1},\ldots,a_{l,N_l}\}$ denote all activation values collected from the calibration set. MinMax calibration uses $\alpha_l^{\mathrm{MM}}=\max_{a\in\mathcal{A}_l}a$, while the fixed percentile baseline uses $\alpha_l^{\mathrm{P99.9}}=\operatorname{Percentile}_{99.9}(\mathcal{A}_l)$. The same 99.9 percentile level is used at all activation sites, although each threshold is computed from the calibration activations of the corresponding site.

Layer-wise MSE-Selected Activation Clipping selects a threshold separately for each post-ReLU activation site. The search is performed over the finite percentile set
\begin{equation}
\mathcal{P}=\{99.0,99.5,99.9,99.95,100.0\}.
\label{eq:candidate_set}
\end{equation}
For each $p\in\mathcal{P}$, the candidate threshold is $\alpha_l(p)=\operatorname{Percentile}_{p}(\mathcal{A}_l)$. The selected percentile is
\begin{equation}
\begin{aligned}
p_l^{\ast} &= \arg\min_{p\in\mathcal{P}} E_l(p), \\
E_l(p) &= \frac{1}{N_l}\sum_{i=1}^{N_l}\left(a_{l,i}-\tilde{a}_{l,i}(p)\right)^2.
\end{aligned}
\label{eq:selected_percentile}
\end{equation}
and the final threshold is $\alpha_l^{\ast}=\alpha_l(p_l^{\ast})$. This procedure is a discrete search over the predefined candidate set, not a continuous global optimization of the clipping threshold. After calibration, the selected activation quantization parameters are attached to the corresponding post-ReLU modules and evaluated in the simulated PTQ pipeline.

\begin{algorithm}[t]
\caption{Layer-wise MSE-Selected Activation Clipping}
\label{alg:mse_selected_clipping}
\begin{algorithmic}[1]
\REQUIRE Trained FP32 model, calibration set, post-ReLU activation sites, percentile set $\mathcal{P}$
\ENSURE Selected thresholds $\{\alpha_l^{\ast}\}$ for all observed activation sites
\FOR{each post-ReLU activation site $l$}
    \STATE Collect calibration activations $\mathcal{A}_l$ from the FP32 model
    \FOR{each percentile $p\in\mathcal{P}$}
        \STATE Set $\alpha_l(p)=\operatorname{Percentile}_{p}(\mathcal{A}_l)$
        \STATE Fake-quantize and dequantize $\mathcal{A}_l$ using $\alpha_l(p)$
        \STATE Compute reconstruction error $E_l(p)$ by Eq.~\eqref{eq:selected_percentile}
    \ENDFOR
    \STATE Select $p_l^{\ast}=\arg\min_{p\in\mathcal{P}}E_l(p)$
    \STATE Set $\alpha_l^{\ast}=\alpha_l(p_l^{\ast})$
\ENDFOR
\RETURN $\{\alpha_l^{\ast}\}$
\end{algorithmic}
\end{algorithm}

\subsection{Interpretation of the Local MSE Criterion}
\label{subsec:mse_criterion_interpretation}

The activation reconstruction MSE used in Eq.~\eqref{eq:selected_percentile} is a calibration-time local criterion. It measures how closely the fake-quantized activations match the FP32 activations at a specific observed site. This criterion is useful because it directly captures the rounding-versus-clipping trade-off induced by a candidate threshold. However, it does not include downstream nonlinearities, residual interactions, class margins, or final decision boundaries. Therefore, a lower activation reconstruction MSE is not guaranteed to produce a higher test-set accuracy. The Results and Discussion section evaluates this distinction explicitly.
